\documentclass[
    notheorems,
    11pt,
    compress,
    aspectratio=169
]{beamer}

%\documentclass[a4paper,11pt]{article}
%\usepackage{beamerarticle}

\usepackage{soul}
\usepackage{xcolor}
\sethlcolor{yellow}
\usepackage{tabto}

\mode<presentation>{
    \usetheme{Madrid}
    \usecolortheme{orchid}
    \usefonttheme[onlymath]{serif}
    \setbeamertemplate{blocks}[rounded][shadow=true]
    \setbeamertemplate{navigation symbols}{}
    \setbeamertemplate{enumerate items}[square]
}

\only<article>{
    \usepackage[margin=1in]{geometry}
    \setlength{\baselineskip}{15pt}
}

\only<presentation>{
    \setlength{\baselineskip}{15pt}
}

\usepackage{config}

\graphicspath{{./figures/}}
\bibliography{bibliography.bib}

\tikzstyle{startstop} = [rectangle, rounded corners, minimum width=3cm, minimum height=1cm, text centered, draw=black, fill=red!30]

\tikzstyle{io} = [trapezium, trapezium left angle=70, text width=3cm, trapezium right angle=110, minimum width=3cm, minimum height=1cm, text centered, draw=black, fill=blue!30]

\tikzstyle{process} = [rectangle, minimum width=3cm, text width=4cm, minimum height=1cm, text centered, draw=black, fill=orange!30]

\tikzstyle{decision} = [diamond, minimum width=3cm, minimum height=1cm, text centered, draw=black, fill=green!30]

\tikzstyle{arrow} = [thick, ->, >=stealth]

% -----

\title[L'istruzione di selezione e le condizioni logiche]{Dal problema al programma: le basi della programmazione}
\subtitle{L'istruzione di selezione e le condizioni logiche}
\author{prof. Leonardo Essam Dei Rossi}
\institute[]{ITT "M. Buonarroti" - Trento (TN)}
\date[A.S. 2025/2026]{Anno scolastico 2025/2026}

% -----

\begin{document}

\begin{frame}
    \titlepage
\end{frame}

\only<presentation>{
    \begin{frame}{Indice}
        \tableofcontents
    \end{frame}
}

\only<article>{
    \maketitle
}

% -----

\section{L'istruzione di selezione doppia}

\begin{frame}
    \only<presentation>{
        \frametitle{L'istruzione di selezione doppia (1)}
    }

    Nella descrizione di un \structure{algoritmo} si può verificare che le operazioni da eseguire siano diverse a seconda dei dati inseriti.

    \pause
    \vspace{\baselineskip}

    \begin{block}{Esempio}
        Ad esempio, se dobbiamo andare in vacanza, dobbiamo verificare di non avere “debiti scolastici” a settembre, e quindi si prospettano due soluzioni alternative:

        \begin{itemize}
            \item \alert{Se} abbiamo tutte sufficienze, \alert{allora} andiamo in vacanza;
            \item \alert{Altrimenti} dobbiamo studiare tutta l’estate.
        \end{itemize}
    \end{block}
\end{frame}

\begin{frame}
    \only<presentation>{
        \frametitle{L'istruzione di selezione doppia (2)}
    }

    Nel diagramma a blocchi questa situazione è illustrata ricorrendo a un elemento grafico specifico, detto \structure{blocco di decisione} o \structure{blocco condizionale}:

    \pause

    \begin{figure}[htbp]
        \centering

        \begin{tikzpicture}[scale=0.85, transform shape]
            \node (co1) [decision] {Condizione};

            \pause

            \node (pr1) [process, below left of=co1, yshift=-1.5cm, xshift=-2cm] {Elaborazione 1};
            \draw [arrow] (co1) -| node[anchor=east] {Yes (VERO)} (pr1);

            \pause
            
            \node (pr2) [process, below right of=co1, yshift=-1.5cm, xshift=2cm] {Elaborazione 2};
            \draw [arrow] (co1) -| node[anchor=west] {No (FALSO)} (pr2);
        \end{tikzpicture}
        
        \caption{Blocco di decisione (condizionale) ramificato}
        \label{fig:11}
    \end{figure}
\end{frame}

\begin{frame}
    \only<presentation>{
        \frametitle{L'istruzione di selezione doppia (3)}
    }

    Di seguito alcuni esempi di istruzioni di confronto che possiamo scrivere nei blocchi condizionali:

    \vspace{\baselineskip}

    \begin{columns}
        \column{0.5\textwidth}
            \begin{itemize}
                \item Ho abbastanza soldi comprare la moto?
                \item Il numero 50 è maggiore di 0?
                \item Mario è un alunno di questa classe?
            \end{itemize}

        \column{0.5\textwidth}
            \begin{itemize}
                \item Il serbatoio è pieno?
                \item Giochi a calcio?
                \item 27 è un numero pari?
            \end{itemize}
    \end{columns}
\end{frame}

\begin{frame}
    \only<presentation>{
        \frametitle{L'istruzione di selezione doppia (4)}
    }

    \begin{block}{}
        A queste domande si può solo rispondere \alert{SÌ} oppure \alert{NO}, esistono cioè soltanto \structure{due alternative}: si tratta dunque di “istruzioni di confronto” che prendono il nome di \alert{condizioni logiche} dalle quali, a seconda del valore della risposta, è sempre necessario intraprendere \alert{percorsi alternativi}, cioè in base alla risposta è necessario eseguire operazioni diverse.
    \end{block}
\end{frame}

\begin{frame}
    \only<presentation>{
        \frametitle{L'istruzione di selezione doppia (5)}
    }
    
    Li rappresentiamo graficamente nei diagrammi a blocchi con uno schema come il seguente:

    \begin{columns}
        \column{0.5\textwidth}
            \begin{figure}[htbp]
                \centering
        
                \begin{tikzpicture}[scale=0.75, transform shape]
                    \node (co1) [decision, text width=2cm] {Sei stato promosso?};
        
                    \node (pr1) [process, below left of=co1, yshift=-1.5cm, xshift=-2cm] {Vai alla classe successiva};
                    \node (pr2) [process, below right of=co1, yshift=-1.5cm, xshift=2cm] {Ripeti la stessa classe};
        
                    \draw [arrow] (co1) -| node[anchor=east] {SÌ} (pr1);
                    \draw [arrow] (co1) -| node[anchor=west] {NO} (pr2);
                \end{tikzpicture}
            \end{figure}

        \pause

        \column{0.5\textwidth}
            \begin{figure}[htbp]
                \centering
        
                \begin{tikzpicture}[scale=0.75, transform shape]
                    \node (co1) [decision, text width=2cm] {Hai 18 anni?};
        
                    \node (pr1) [process, below left of=co1, yshift=-1.5cm, xshift=-2cm] {Puoi prendere la patente};
                    \node (pr2) [process, below right of=co1, yshift=-1.5cm, xshift=2cm] {Non puoi guidare};
        
                    \draw [arrow] (co1) -| node[anchor=east] {SÌ} (pr1);
                    \draw [arrow] (co1) -| node[anchor=west] {NO} (pr2);
                \end{tikzpicture}
            \end{figure}
    \end{columns}
\end{frame}

\begin{frame}
    \only<presentation>{
        \frametitle{L'istruzione di selezione doppia (6)}
    }

    \begin{columns}
        \column{0.45\textwidth}
            \begin{block}{}
                Possiamo osservare che, a seconda della risposta data alla domanda (cioè al \alert{test}), le strade “si dividono”, vengono eseguite operazioni diverse nei due “rami” ma, successivamente, si ricongiungono.
            \end{block}

            \pause

            \only<presentation>{
                \hl{La presenza di istruzioni in entrambi i rami del diagramma fa sì che in questo caso si parli di \textbf{selezione doppia}.}

                \pause
            }

            \only<article>{
                La presenza di istruzioni in entrambi i rami del diagramma fa sì che in questo caso si parli di \textbf{selezione doppia}.
            }

            \only<presentation>{
                \vspace{\baselineskip}
                
                Scriviamo il diagramma a blocchi per l’algoritmo “semplificato” che descrive la procedura per l’acquisto di una moto.
            }

        \column{0.45\textwidth}
            \only<presentation>{
                \pause
                \begin{figure}
                    \centering
    
                    \begin{tikzpicture}[scale=0.75, transform shape]
                        \node (start) [startstop] {Inizio};
    
                        \node (io1) [io, below of=start, yshift=-0.5cm] {Leggi: costo della moto};
                        \draw [arrow] (start) -- (io1);
    
                        \node (io2) [io, below of=io1, yshift=-0.5cm] {Leggi: soldi a disposizione};
                        \draw [arrow] (io1) -- (io2);
                    \end{tikzpicture}
                \end{figure}
    
                \begin{center}
                    \textit{(continua)}
                \end{center}
            }
    \end{columns}
\end{frame}

\begin{frame}
    \only<presentation>{
        \frametitle{L'istruzione di selezione doppia (7)}
    }

    \begin{columns}
        \column{0.45\textwidth}
            \only<presentation>{
                \begin{block}{}
                    Possiamo osservare che, a seconda della risposta data alla domanda (cioè al \alert{test}), le strade “si dividono”, vengono eseguite operazioni diverse nei due “rami” ma, successivamente, si ricongiungono.
                \end{block}
    
                \hl{La presenza di istruzioni in entrambi i rami del diagramma fa sì che in questo caso si parli di \textbf{selezione doppia}.}
            }

            \vspace{\baselineskip}

            \only<presentation>{
                Scriviamo il diagramma a blocchi per l’algoritmo “semplificato” che descrive la procedura per l’acquisto di una moto.
            }

            \only<article>{
                Scriviamo il diagramma a blocchi per l’algoritmo “semplificato” che descrive la procedura per l’acquisto di una moto:
            }

        \column{0.45\textwidth}
            \only<presentation>{
                \begin{figure}
                    \centering
    
                    \begin{tikzpicture}[scale=0.75, transform shape]
                        \node (co1) [decision, text width=2cm] {I soldi che hai sono sufficienti?};
            
                        \node (pr1) [process, below left of=co1, yshift=-1.5cm, xshift=-2cm] {Stampa: "puoi comprare la moto"};
                        \node (pr2) [process, below right of=co1, yshift=-1.5cm, xshift=2cm] {Stampa: "devi guadagnare altri soldi!"};
            
                        \draw [arrow] (co1) -| node[anchor=east] {SÌ} (pr1);
                        \draw [arrow] (co1) -| node[anchor=west] {NO} (pr2);
    
                        \node (stop) [startstop, below of=pr1, yshift=-1cm, xshift=2.7cm] {Fine};
    
                        \draw [arrow] (pr1) -| node[anchor=north] {} (stop);
                        \draw [arrow] (pr2) -| node[anchor=north] {} (stop);
                    \end{tikzpicture}
                \end{figure}
            }
    \end{columns}
\end{frame}

\begin{frame}
    \only<presentation>{
        \frametitle{L'istruzione di selezione doppia (8)}
    }

    \begin{figure}[htbp]
        \centering

        \begin{tikzpicture}[scale=0.75, transform shape]
            \node (start) [startstop] {Inizio};

            \node (io1) [io, below of=start, yshift=-0.5cm] {Leggi: costo della moto};
            \draw [arrow] (start) -- (io1);

            \node (io2) [io, below of=io1, yshift=-0.5cm] {Leggi: soldi a disposizione};
            \draw [arrow] (io1) -- (io2);


            \node (co1) [decision, right of=start, text width=2cm, xshift=8cm] {I soldi che hai sono sufficienti?};
        
            \node (pr1) [process, below left of=co1, yshift=-1.5cm, xshift=-2cm] {Stampa: "puoi comprare la moto"};
            \node (pr2) [process, below right of=co1, yshift=-1.5cm, xshift=2cm] {Stampa: "devi guadagnare altri soldi!"};

            \draw [arrow] (co1) -| node[anchor=east] {SÌ} (pr1);
            \draw [arrow] (co1) -| node[anchor=west] {NO} (pr2);

            \node (stop) [startstop, below of=pr1, yshift=-1cm, xshift=2.7cm] {Fine};

            \draw [arrow] (pr1) -| node[anchor=north] {} (stop);
            \draw [arrow] (pr2) -| node[anchor=north] {} (stop);

            \draw [arrow] (io2.south) |- ++(3,-0.5) |- ([yshift=0.6cm]co1.north) -- (co1.north);
        \end{tikzpicture}

        \caption{Diagramma di flusso completo}
        \label{fig:12}
    \end{figure}
\end{frame}

% -----

\only<article>{
    \newpage
}

\section{Il linguaggio di progetto}

\begin{frame}[fragile]
    \only<presentation>{
        \frametitle{Il linguaggio di progetto (1)}
    }

    Iniziamo a definire un "nostro" \structure{linguaggio di progetto}, ovvero un insieme di convenzioni da rispettare quando si scrive un algoritmo in pseudocodifica.

    \vspace{\baselineskip}

    \only<article>{
        \noindent
        Definiamo la struttura base di un programma:
        
        \vspace{\baselineskip}

        \noindent
        \fbox{\parbox{\textwidth}{
            \structure{\texttt{Inizio}}
            \\
            ...
            \\
            \structure{\texttt{Fine}}
        }}
    }

    \only<presentation>{
        Definiamo la struttura base di un programma:

        \vspace{\baselineskip}
    
        \fbox{\parbox{0.95\textwidth}{
            \structure{\texttt{Inizio}}
            \\
            ...
            \\
            \structure{\texttt{Fine}}
        }}
    }
\end{frame}

\begin{frame}[fragile]
    \only<presentation>{
        \frametitle{Il linguaggio di progetto (2)}
    }

    \only<presentation>{
        Definiamo ora le istruzioni di input e output:

        \vspace{\baselineskip}
    
        \fbox{\parbox{0.95\textwidth}{
            \structure{\texttt{Inizio}}
            \\
            \structure{\texttt{Leggi (distanza\_km):}} "Inserisci la distanza in km"
            \\
            ...
            \\
            \structure{\texttt{Stampa (distanza\_km):}} "La distanza in km è: [1]"
            \\
            \structure{\texttt{Fine}}
        }}
    }

    \only<article>{
        \vspace{\baselineskip}
        \noindent
        Definiamo ora le istruzioni di input e output:

        \vspace{\baselineskip}

        \noindent
        \fbox{\parbox{\textwidth}{
            \structure{\texttt{Inizio}}
            \\
            \structure{\texttt{Leggi (distanza\_km):}} "Inserisci la distanza in km"
            \\
            ...
            \\
            \structure{\texttt{Stampa (distanza\_km):}} "La distanza in km è: [1]"
            \\
            \structure{\texttt{Fine}}
        }}
    }
\end{frame}

\begin{frame}
    \only<presentation>{
        \frametitle{Il linguaggio di progetto (3)}
    }

    \only<article>{
        \vspace{\baselineskip}
        \noindent
        E l'istruzione di elaborazione:

        \vspace{\baselineskip}

        \noindent
        \fbox{\parbox{\textwidth}{
            \structure{\texttt{Inizio}}
            \\
            \structure{\texttt{Leggi (distanza\_km):}} "Inserisci la distanza in km"
            \\
            \structure{\texttt{Calcola:}} (\texttt{distanza\_m = distanza\_km * 1000})
            \\
            \structure{\texttt{Stampa (distanza\_km, distanza\_m):}} "La distanza in km è [1], mentre in metri è [2]"
            \\
            \structure{\texttt{Fine}}
        }}
    }

    \only<presentation>{
        E l'istruzione di elaborazione:
    
        \vspace{\baselineskip}
    
        \fbox{\parbox{0.95\textwidth}{
            \structure{\texttt{Inizio}}
            \\
            \structure{\texttt{Leggi (distanza\_km):}} "Inserisci la distanza in km"
            \\
            \structure{\texttt{Calcola:}} (\texttt{distanza\_m = distanza\_km * 1000})
            \\
            \structure{\texttt{Stampa (distanza\_km, distanza\_m):}} "La distanza in km è [1], mentre in metri è [2]"
            \\
            \structure{\texttt{Fine}}
        }}
    }
\end{frame}

\begin{frame}
    \only<presentation>{
        \frametitle{Il linguaggio di progetto (5)}
    }

    \only<article>{
        \vspace{\baselineskip}
        \noindent
        Per finire, l'istruzione di decisione:

        \vspace{\baselineskip}

        \noindent
        \fbox{\parbox{\textwidth}{
            \structure{\texttt{Inizio}}
            \\
            \structure{\texttt{Leggi (distanza\_km):}} "Inserisci la distanza in km"
            \\
            \structure{\texttt{Calcola:}} (\texttt{distanza\_m = distanza\_km * 1000})
            \\
            \\
            \structure{\texttt{Se (distanza\_m > 1000), allora:}}
            \\
            \structure{-~~~~\texttt{Stampa ():}} "La distanza è più di 1000 metri!"
            \\
            \structure{\texttt{Altrimenti:}}
            \\
            \structure{-~~~~\texttt{Stampa ():}} "La distanza è meno di 1000 metri!"
            \\
            \\
            \structure{\texttt{Fine}}
        }}
    }

    \only<presentation>{
        Per finire, l'istruzione di decisione:
    
        \vspace{\baselineskip}
    
        \fbox{\parbox{0.95\textwidth}{
            \structure{\texttt{Inizio}}
            \\
            \structure{\texttt{Leggi (distanza\_km):}} "Inserisci la distanza in km"
            \\
            \structure{\texttt{Calcola:}} (\texttt{distanza\_m = distanza\_km * 1000})
            \\
            \\
            \structure{\texttt{Se (distanza\_m > 1000), allora:}}
            \\
            \structure{-~~~~\texttt{Stampa ():}} "La distanza è più di 1000 metri!"
            \\
            \structure{\texttt{Altrimenti:}}
            \\
            \structure{-~~~~\texttt{Stampa ():}} "La distanza è meno di 1000 metri!"
            \\
            \\
            \structure{\texttt{Fine}}
        }}
    }
\end{frame}

% -----

\only<presentation>{
    \subsection{Esercizio \#1}
}

\only<article>{
    \subsection{Esercizio \#1 (Tratto da: \cite[p. 350]{ISBN_8820378434})}
}

\begin{frame}
    \only<presentation>{
        \frametitle{Esercizio \#1}

        \textbf{Si consideri la \structure{Figura \ref{fig:12}}.}
        \\
        Scrivere la pseudocodifica corrispondente in linguaggio formale.
    }
    

    \only<presentation>{
        \vfill
    }

    \only<article>{
        Si consideri la \structure{Figura \ref{fig:12}}.
        \\
        Scrivere la pseudocodifica corrispondente in linguaggio formale.
    
        \vspace{\baselineskip}
    }
\end{frame}

% ----

\only<presentation>{
    \subsection{Esercizio \#2}
}

\only<article>{
    \subsection{Esercizio \#2 (Tratto da: \cite[p. 355]{ISBN_8820378434})}
}

\begin{frame}
    \only<presentation>{
        \frametitle{Esercizio \#2}
    }

    Si scriva il diagramma di flusso e la relativa pseudocodifica in linguaggio formale del seguente esercizio: "\textit{leggi due numeri e visualizza sullo schermo il maggiore di essi}".

    \only<presentation>{
        \vfill
    }

    \only<article>{
        \vspace{\baselineskip}
    }
\end{frame}

% -----

\only<article>{
    \newpage
}

\section{La selezione semplice}

\begin{frame}
    \only<presentation>{
        \frametitle{La selezione semplice (1)}
    }

    È possibile che le operazioni da eseguire siano presenti in un solo ramo, quando ad esempio si devono effettuare operazione solo in un caso dell’esito del test.

    \vspace{\baselineskip}

    In queste situazioni si parla di \structure{selezione semplice}.
\end{frame}

\begin{frame}
    \only<presentation>{
        \frametitle{La selezione semplice (2)}
    }

    \only<article>{
        \vspace{\baselineskip}
    }

    \begin{block}{Esempio}
        Scriviamo un programma che trasforma i numeri negativi in positivi: nel caso che il numero inserito sia già positivo non dobbiamo effettuare nessuna operazione mentre se è negativo lo moltiplichiamo per –1, in modo da cambiarne il segno.
    \end{block}

    \vspace{\baselineskip}

    \pause

    \only<article>{
        \noindent
        Una possibile soluzione è la seguente:

        \vspace{\baselineskip}

        \noindent
        \fbox{\parbox{\textwidth}{
            \structure{\texttt{Inizio}}
            \\
            \structure{\texttt{Leggi (numero):}} "Inserisci un numero"
            \\
            \structure{\texttt{Se (numero < 0), allora:}}
            \\
            \structure{-~~~~\texttt{Calcola:}} (\texttt{numero = numero * (-1)})
            \\
            \structure{\texttt{Stampa (numero):}} "Il numero è [1]"
            \\
            \structure{\texttt{Fine}}
        }}
    }

    \only<presentation>{
        Una possibile soluzione è la seguente:

        % \vspace{\baselineskip}
    
        \fbox{\parbox{0.95\textwidth}{
            \structure{\texttt{Inizio}}
            \\
            \structure{\texttt{Leggi (numero):}} "Inserisci un numero"
            \\
            \structure{\texttt{Se (numero < 0), allora:}}
            \\
            \structure{-~~~~\texttt{Calcola:}} (\texttt{numero = numero * (-1)})
            \\
            \structure{\texttt{Stampa (numero):}} "Il numero è [1]"
            \\
            \structure{\texttt{Fine}}
        }}
    }
\end{frame}

% -----

\subsection{Esercizio \#3}

\begin{frame}
    \only<presentation>{
        \frametitle{Esercizio \#3}
    }

    \textbf{Si consideri la pseudocodifica formale presentata.}
    \\
    Scrivere il corrispettivo diagramma di flusso.
\end{frame}

% -----

\subsection{Istruzione di selezione semplice vs. doppia}

\begin{frame}
    \only<presentation>{
        \frametitle{Istruzione di selezione semplice vs. doppia}
    
        \begin{block}{Nota bene}
            L’istruzione di \structure{selezione semplice} può quindi essere considerata un caso particolare di \structure{selezione doppia} dove viene a mancare il ramo \alert{altrimenti}; non è invece possibile avere un’istruzione semplice senza il ramo \alert{allora}, con cioè il solo ramo \alert{altrimenti}.
        \end{block}
    }

    \only<article>{
        \textbf{Nota bene.} L’istruzione di \structure{selezione semplice} può quindi essere considerata un caso particolare di \structure{selezione doppia} dove viene a mancare il ramo \alert{altrimenti}; non è invece possibile avere un’istruzione semplice senza il ramo \alert{allora}, con cioè il solo ramo \alert{altrimenti}.
    }
\end{frame}

% -----

\section{Esercizi}

\subsection{Esercizio \#4}

\begin{frame}
    \only<presentation>{
        \frametitle{Esercizio \#4}
    }

    Si scrivano il diagramma di flusso e la pseudocodifica in linguaggio formale di un programma che, dati in input tre numeri, stabilisca se tali numeri siano o meno una terna pitagorica.

    \vspace{\baselineskip}

    \only<presentation>{
        \begin{block}{Definizione: Terna pitagorica}
            Tre numeri formano una terna pitagorica se il quadrato del primo è uguale alla somma dei quadrati del secondo e del terzo.
    
            \vspace{\baselineskip}
    
            \textbf{Esempio:} \((3, 4, 5) \to (9, 16, 25) \Rightarrow 9 + 16 = 25\)
        \end{block}
    }

    \only<article>{
        \textbf{Definizione.} Tre numeri formano una terna pitagorica se il quadrato del primo è uguale alla somma dei quadrati del secondo e del terzo.

        \vspace{\baselineskip}

        \textbf{Esempio.} \((3, 4, 5) \to (9, 16, 25) \Rightarrow 9 + 16 = 25\)
    }

    \vfill

    \only<presentation>{
        Tratto da: \structure{\cite[p. 355]{ISBN_8820378434}}
    }
\end{frame}

% -----

\only<article>{
    \newpage
}

\only<presentation>{
    \subsection{Esercizio \#5}
}

\only<article>{
    \subsection{Esercizio \#5 (Tratto da: \cite[p. 358]{ISBN_8820378434})}
}

\begin{frame}
    \only<presentation>{
        \frametitle{Esercizio \#5}
    }

    Si consideri la seguente pseudocodifica formale:

    \vspace{\baselineskip}

    \only<presentation>{
        \fbox{\parbox{0.95\textwidth}{
            \structure{\texttt{Inizio}}
            \\
            \structure{\texttt{Leggi (numero1, numero2):}} "Inserisci due numeri"
            \\
            \structure{\texttt{Se (numero1 < numero2), allora:}}
            \\
            \structure{-~~~~\texttt{Stampa (numero2):}} "[1]"
            \\
            \structure{\texttt{Leggi (numero3):}} "Inserisci un numero"
            \\
            \structure{\texttt{Se (numero2 < numero3), allora:}}
            \\
            \structure{-~~~~\texttt{Stampa (numero2):}} "[1]"
            \\
            \structure{\texttt{Altrimenti:}}
            \\
            \structure{-~~~~\texttt{Stampa (numero3):}} "[1]"
            \\
            \structure{\texttt{Fine}}
        }}
    }

    \only<article>{
        \noindent
        \fbox{\parbox{\textwidth}{
            \structure{\texttt{Inizio}}
            \\
            \structure{\texttt{Leggi (numero1, numero2):}} "Inserisci due numeri"
            \\
            \structure{\texttt{Se (numero1 < numero2), allora:}}
            \\
            \structure{-~~~~\texttt{Stampa (numero2):}} "[1]"
            \\
            \structure{\texttt{Leggi (numero3):}} "Inserisci un numero"
            \\
            \structure{\texttt{Se (numero2 < numero3), allora:}}
            \\
            \structure{-~~~~\texttt{Stampa (numero2):}} "[1]"
            \\
            \structure{\texttt{Altrimenti:}}
            \\
            \structure{-~~~~\texttt{Stampa (numero3):}} "[1]"
            \\
            \structure{\texttt{Fine}}
        }}
    }

    \vspace{\baselineskip}

    Cosa viene visualizzato inserendo i valori \((5, 2, 7)\)?

    \only<presentation>{
        \vfill
    }

    \only<article>{
        \vspace{\baselineskip}
    }
\end{frame}

% -----

\begin{frame}
    \only<presentation>{
        \frametitle{Riferimenti e approfondimenti}    
    }
    
    \printbibliography
\end{frame}

% -----

\end{document}
